\section{Relations I}
\begin{questions}
    % Relation Discriminator
    \item \begin{partquestions}{\alph*}
        \item Since $\emptyset \subseteq A\times B$ and $\emptyset \subseteq A^2$, thus $R_1$ is a relation from $A$ to $B$ as well as a relation on $A$.
        \item Note $(3, 3) \notin A\times B$, so $R_2$ is not a relation from $A$ to $B$. However $R_2 \subseteq A^2$ so $R_2$ is a relation on $A$.
        \item Observe $R_3 \subseteq A\times B$ so $R_3$ is a relation from $A$ to $B$. But $(4, 8) \notin A^2$ so $R_3$ is not a relation on $A$.
        \item Note $2 \notin A \times B$ and $2 \notin A^2$, so $R_4$ is neither a relation from $A$ to $B$ nor a relation on $A$.
        \item One sees $(6, 8) \notin A\times B$ and $(6, 8) \notin A^2$ and thus $R_5$ is neither a relation from $A$ to $B$ nor a relation on $A$.
        \item Note that $(5, 5) \notin A \times B$ so $R_6$ is not a relation from $A$ to $B$. However clearly $A^2 \subseteq A^2$, which means $R_6$ is a relation on $A$.
        \item Since $(6, 6) \notin A \times B$ and $(6, 6) \notin A^2$, thus $R_7$ is neither a relation from $A$ to $B$ nor a relation on $A$.
    \end{partquestions}

    % Properties of the "Divides" Relation
    \item \begin{partquestions}{\roman*}
        \item We will do this step by step.
        \begin{itemize}
            \item For $x = 2$, notice that $2 \mathrel{D} 4$ (using $m = 2$), $2 \mathrel{D} 6$ (using $m = 3$), and $2 \mathrel{D} 8$ (using $m = 4$).
            \item For $x = 3$, notice that only $3 \mathrel{D} 6$ (using $m = 2$).
            \item For $x = 4$, notice that only $4 \mathrel{D} 8$ (using $m = 2$).
            \item For $x = 5$, $x = 6$, $x = 7$, and $x = 8$, there are no valid $y$ that is related to them.
        \end{itemize}
        Thus $D = \{(2, 4), (2, 6), (2, 8), (3, 6), (4, 8)\}$. Hence we see
        \begin{align*}
            D^{-1} &= \{(4, 2), (6, 2), (8, 2), (6, 3), (8, 4)\} & (\text{Swap } x \text{ and } y)\\
            &= \{(4, 2), (6, 2), (6, 3), (8, 2), (8, 4)\}. & (\text{Rearranging})
        \end{align*}
        
        \item \begin{partquestions}{\alph*}
            \item Domain of $D$ is $\{2, 3, 4\}$.
            \item Codomain of $D$ is $S = \{2,3,4,5,6,7,8\}$.
            \item Range of $D$ is $\{4, 6, 8\}$.
        \end{partquestions}
    \end{partquestions}

    % Inverse Relation of Inverse Relation
    \item Note that $\left(R^{-1}\right)^{-1} \subseteq A \times B$. One sees
    \begin{align*}
        &(a, b) \in \left(R^{-1}\right)^{-1}\\
        \iff& (b, a) \in R^{-1} & (\text{Definition of inverse relation})\\
        \iff& (a, b) \in R & (\text{Definition of inverse relation})
    \end{align*}
    Thus any element in $\left(R^{-1}\right)^{-1}$ is also in $R$ and vice versa. Hence $\left(R^{-1}\right)^{-1} = R$.
    
    % Properties of Compositions
    \item \begin{partquestions}{\alph*}
        \item True. Let $A = \{1\}$, $B = \{2, 3\}$, $C = \{4\}$, $R = \{(1,2)\}$, and $S = \{(3,4)\}$. Thus $S \circ R = \emptyset$ and so the domain of $S \circ R$ is empty.
        \item False, see \textbf{(a)}.
        \item True. Let $A = \{1\}$, $B = \{2\}$, $C = \{3\}$, $R = \{(1,2)\}$, and $S = \{(2,3)\}$. Thus $S \circ R = \{(1,3)\}$ and so the domain of $S \circ R$ is $\{1\} = A$.
        \item False. The codomain of $S \circ R$ is $C$ which is given to be non-empty.
        \item True. This is the definition of the codomain of a relation.
        \item False. If the codomain is indeed $B$, that means $C = B$ (since $C$ is the actual codomain). But $B$ and $C$ are distinct, so this cannot happen.
        \item True, see \textbf{(a)}.
        \item False, see \textbf{(a)}.
        \item True, see \textbf{(c)}.
    \end{partquestions}

    % Composition Commutativity?
    \item \begin{partquestions}{\alph*}
        \item False. Let $A = \{1\}$ and $B = \{2\}$. Let $R = \{(1, 2)\}$ and $S = \{(2, 1)\}$. We see $S \circ R = \{(1, 1)\} \neq \{(2, 2)\} = R \circ S$. Thus it is not always the case that $S \circ R = R \circ S$.
        \item False. See \textbf{(c)}.
        \item True. Let $A = B = \{1, 2\}$. Let $R = S = \{(1, 2), (2, 1)\}$. Thus $S \circ R = \{(1, 1), (2, 2)\} = R \circ S$. So it is possible for $S \circ R$ and $R \circ S$ to be the same.
    \end{partquestions}

    % Shoes and Socks
    \item We work step by step.
    \begin{align*}
        &(c, a) \in (S \circ R)^{-1}\\
        \iff& (a, c) \in S \circ R & (\text{Definition of inverse relation})\\
        \iff& \exists b \in B \text{ s.t. } (a \mathrel{R} b \land b \mathrel{S} c) & (\text{Definition of composition})\\
        \iff& \exists b \in B \text{ s.t. } (b \mathrel{R^{-1}} a \land c \mathrel{S^{-1}} b) & (\text{Definition of inverse relation})\\
        \iff& \exists b \in B \text{ s.t. } (c \mathrel{S^{-1}} b \land b \mathrel{R^{-1}} a) & (\text{Commutativity of } \land)\\
        \iff& (c, a) \in R^{-1} \circ S^{-1} & (\text{Definition of composition})
    \end{align*}
    Thus any element in $(S \circ R)^{-1}$ is also in $R^{-1} \circ S^{-1}$ and vice versa, thereby meaning that $(S \circ R)^{-1} = R^{-1} \circ S^{-1}$.

    % Composition of Relation and its Inverse
    \item The statement is false. Consider $A = B = \{1, 2, 3\}$ and the relation $R = \{(1, 3), (2, 3)\}$. We see that $R^{-1} = \{(3, 1), (3, 2)\}$. However, note that $(1, 2) \in R^{-1} \circ R$ since $1 \mathrel{R} 3$ and $3 \mathrel{R^{-1}} 2$. Thus $R^{-1} \circ R \not\subseteq \{(x, y) \in A^2 \vert x = y\} = \{(1, 1), (2, 2), (3, 3)\}$.
\end{questions}
